\chapter{Peripheral Angioplasty and Catheter Ablation}

\section*{Introduction \& Clinical Indications}
This chapter covers two important catheter-based intervention families: endovascular angioplasty or stenting for selected peripheral arterial disease, and electrophysiology catheter ablation for selected tachyarrhythmias. They are not interchangeable operations and require different specialist teams, imaging, anticoagulation, and follow-up. Each is included because modern procedural care extends beyond open surgery.

Peripheral angioplasty is considered for lifestyle-limiting claudication despite medical therapy and exercise, or for chronic limb-threatening ischemia when revascularization can relieve ischemia and support wound healing. Acute limb ischemia requires urgent vascular assessment and may require thrombectomy, thrombolysis, open bypass, or endovascular treatment. Catheter ablation is used for symptomatic or clinically important arrhythmias such as supraventricular tachycardia, atrial flutter, selected atrial fibrillation, and some ventricular arrhythmias.

\section*{Relevant Surgical Anatomy}
Peripheral procedures may involve the aortoiliac, femoropopliteal, or infrapopliteal circulation; lesion length, calcification, runoff, inflow, and collateral vessels determine strategy. Ablation targets are defined by the cardiac conduction system and arrhythmia mechanism. The atrioventricular node, His bundle, coronary arteries, phrenic nerves, pulmonary veins, and esophagus may be near ablation targets.

\section*{Preoperative Workup \& Preparation}
For peripheral disease, document pulses, skin, wounds, ankle-brachial or toe pressures, duplex or cross-sectional imaging, renal function, diabetes, smoking, and walking limitation. For ablation, obtain ECG documentation, echocardiography when indicated, rhythm monitoring, medication review, thromboembolic-risk assessment, and imaging for selected left-atrial procedures. Discuss contrast injury, bleeding, embolization, vessel perforation, restenosis, stroke, arrhythmia, heart block, tamponade, pulmonary-vein stenosis, esophageal injury, and need for repeat treatment.

\section*{Surgical Technique \& Procedure}
Peripheral angioplasty uses arterial access, a guidewire across the lesion, balloon dilation, and selective stent or drug-coated technology according to anatomy and evidence. Completion imaging checks flow and distal embolization. Catheter ablation uses intracardiac mapping to identify the arrhythmogenic substrate, followed by radiofrequency or cryothermal energy. Anticoagulation and vascular closure are managed according to the procedure and patient risk.

\section*{Complications \& Risks}
Peripheral risks include access-site hematoma, pseudoaneurysm, dissection, perforation, distal embolization, acute occlusion, contrast-associated kidney injury, restenosis, and limb ischemia. Ablation risks include vascular injury, bleeding, stroke, transient ischemic attack, tamponade, phrenic-nerve injury, pulmonary-vein stenosis, esophageal injury, atrioventricular block, recurrence, and death. Sudden limb pain or pallor, neurologic symptoms, chest pain, severe dyspnea, or uncontrolled bleeding requires emergency assessment.

\section*{Postoperative Care \& Recovery}
Check access sites, distal pulses, renal function, rhythm, neurologic status, and symptoms. Peripheral patients need antiplatelet or antithrombotic therapy as prescribed, supervised exercise, wound care, foot protection, lipid and blood-pressure management, and smoking cessation. Ablation patients need medication and anticoagulation instructions, activity restrictions, rhythm follow-up, and awareness that early transient arrhythmias do not always represent failure.

\section*{Outcomes \& Prognosis}
Endovascular success depends on lesion anatomy, medical therapy, exercise, and surveillance; restenosis and progression remain possible. Ablation can markedly reduce arrhythmia burden and symptoms, but recurrence and repeat procedures occur. Neither procedure removes the need to manage the underlying vascular or cardiac disease.

\section*{References}
\begin{enumerate}
\item American Heart Association/American College of Cardiology. Guidelines for peripheral artery disease and atrial fibrillation.
\item European Society for Vascular Surgery. Guidelines on chronic limb-threatening ischemia and peripheral arterial disease.
\item Heart Rhythm Society/European Heart Rhythm Association. Expert consensus on catheter ablation of atrial fibrillation and other arrhythmias.
\end{enumerate}
